% RoboRT tech report — arXiv-style preprint.
% Compile: pdflatex robort-tech-report.tex
\documentclass[11pt]{article}

\usepackage[margin=1in]{geometry}
\usepackage{graphicx}
\usepackage{booktabs}
\usepackage{amsmath}
\usepackage{url}
\usepackage[hidelinks]{hyperref}

\title{\textbf{RoboRT: A Real-Time C++ Inference Engine for pi0.5 Vision-Language-Action Policies}}
\author{liuwang\\[2pt] Baidu}
\date{}

\begin{document}

\maketitle

\begin{abstract}
Vision-language-action (VLA) models such as pi0.5 have demonstrated strong
performance in robot manipulation, yet their inference latency remains the
dominant bottleneck to real-time closed-loop control. We present RoboRT, a
native C++ inference engine tailored to the flow-matching denoising schedule
of pi0.5. Rather than adapting a generic large-language-model (LLM) runtime,
RoboRT implements a hand-written GGUF reader that loads F32, BF16, and
quantized weights directly, a compute graph that is cached and rebuilt only
when the input shape changes, and an engine-owned noise RNG that renders
every prediction bit-exact reproducible. On an A10 GPU, RoboRT achieves
\textbf{161\,ms/step} (6.2\,Hz), a \textbf{2.3$\times$} speedup over the
LeRobot PyTorch reference at 377\,ms/step (2.7\,Hz) under identical
observation and seed, while retaining 38/40 task success on the LIBERO
benchmark. RoboRT is released under the Apache-2.0 license together with a
deterministic self-check tool that enables bit-exact comparison across
implementations.
\end{abstract}

\section{Introduction}
\label{sec:intro}

Recent vision-language-action (VLA) policies, exemplified by the pi0 family
\cite{pi0}, jointly model vision, language, and action in a single
architecture and have become a de facto baseline for generalist robot
manipulation. A VLA policy must keep pace with the control loop it drives:
higher inference throughput directly translates into higher task success in
dynamic manipulation. The practical target is on the order of 10\,Hz per
inference step; anything slower necessitates aggressive motion smoothing and
degrades closed-loop behavior.

Despite their appeal, deploying these models at interactive rates is
challenging for three reasons. First, \emph{reference implementations are
slow}: the widely used LeRobot \cite{lerobot} stack runs pi0.5 at
377\,ms/step on an A10, sustaining only 2.7\,Hz---a small fraction of the
10\,Hz target---and leaving little headroom for a reactive controller.
Second, \emph{generic LLM runtimes are ill-suited to the workload}: mature
inference engines are optimized for autoregressive token generation with
continuous batching, speculative decoding, and KV-cache paging, whereas
pi0.5 predicts action chunks through an iterative flow-matching denoising
loop \cite{flowmatching} in which the entire chunk is denoised in parallel
over a fixed number of steps, all sharing a single precomputed
visual-language prefix. Grafting the model onto such a runtime introduces
substantial unused machinery and a scheduling model that is fundamentally
mismatched to the denoising workload. Third, \emph{the deployment chain is
fragmented}: vision encoding, prefix caching, noise scheduling, and action
denormalization are usually assembled from separate, independently
maintained implementations, rendering results difficult to reproduce and
still more difficult to compare across codebases.

We address these limitations with \textbf{RoboRT}, a self-contained C++
inference engine for pi0.5. The engine is native rather than an adapter over
an LLM framework: it reads weights with a hand-written GGUF reader, executes
a compute graph that is cached across steps and rebuilt only when the input
shapes change, and maintains the noise RNG as engine state, rendering every
prediction bit-exact reproducible from a fixed seed. RoboRT builds solely on
generic compute primitives---the ggml graph runtime with its CPU/CUDA
backends and the SigLIP \cite{siglip} vision tower---and contains no
LLM-specific layers.

At 161\,ms/step (6.2\,Hz), RoboRT closes more than half of the gap toward
interactive control rates on the same A10 hardware, and its per-phase
breakdown shows that denoising---not vision encoding---dominates,
indicating clear headroom for further optimization
(Sec.~\ref{sec:exp-latency}).

Our contributions are:
\begin{itemize}
  \item \textbf{A native pi0.5 inference engine} (Sec.~\ref{sec:method}):
        a hand-written GGUF reader, a cached compute graph, and a
        flow-matching denoising loop, with no dependence on an LLM runtime.
  \item \textbf{A reproducibility mechanism} (Sec.~\ref{sec:repro}):
        deterministic engine-state RNG semantics and a self-check tool that
        emits 1,600 floats bit-exactly for cross-implementation regression.
  \item \textbf{Measured results} (Sec.~\ref{sec:experiments}):
        161\,ms/step on an A10 (2.3$\times$ over LeRobot) and 38/40 task
        success on LIBERO with seed 42, consistent with the reference policy.
\end{itemize}

\section{Related Work}
\label{sec:related}

\textbf{pi0 and pi0.5.} pi0 \cite{pi0} is an open-source VLA built on the
PaliGemma vision-language backbone with a flow-matching action expert;
pi0.5 \cite{pi05} extends it with a stronger backbone while retaining the
same denoising interface. Both predict a chunk of actions (the suffix) by
iteratively denoising Gaussian noise conditioned on vision, language, and
proprioceptive state.

\textbf{VLA inference stacks.} LeRobot \cite{lerobot} provides the canonical
PyTorch implementation of pi0 and pi0.5 and serves as the reference against
which we compare. Its latency is dominated by Python overhead and eager
execution. Serving-oriented runtimes (e.g., SGLang \cite{sglang}) offer
request batching and continuous KV-cache management, but they target
autoregressive LLM serving rather than the non-autoregressive denoising loop
of action experts.

\textbf{Flow matching and diffusion inference.} Flow matching
\cite{flowmatching} and diffusion models \cite{ddpm} share an iterative
refinement schedule. Production inference for these models optimizes the
number of steps, the schedule, and the per-step compute---the same design
knobs that RoboRT exposes---rather than token-level batching.

\textbf{Graph runtimes and vision encoders.} ggml \cite{ggml} provides a
dependency-free computation-graph runtime with CPU and CUDA backends. We use
only its generic operators and implement all pi0.5-specific graph
construction ourselves. The vision tower is the SigLIP \cite{siglip}
encoder, executed as a standalone graph whose output is incorporated into
the language-model prefix.

\textbf{Benchmarks.} LIBERO \cite{libero} is a suite of scripted
long-horizon manipulation tasks that has become a standard benchmark for
evaluating VLA policies; we use it to confirm that RoboRT preserves the
behavior of the reference policy end to end.

\section{Method}
\label{sec:method}

\begin{figure}[t]
  \centering
  \fbox{\parbox{0.92\linewidth}{\centering
  \small
  observation (image, language tokens, state) \\
  $\downarrow$ \\
  SigLIP vision encoding $\rightarrow$ prefix: vision+language embeddings (KV cached) \\
  $\downarrow$ \\
  flow-matching denoising loop: $\mathbf{x}_{T} \xrightarrow{~~\epsilon_\theta~~} \mathbf{x}_{0}$ \\
  $\downarrow$ \\
  denormalization $\rightarrow$ action chunk ($n_{\text{steps}} \times d_{\text{act}}$)
  }}
  \caption{End-to-end data flow of RoboRT for one pi0.5 predict.}
  \label{fig:pipeline}
\end{figure}

Figure~\ref{fig:pipeline} illustrates the end-to-end data flow. The engine
is organized as a single library (\texttt{vla-pi05}) with a public C++ API:
a \texttt{Config} describing the model dimensions, an \texttt{Inputs} struct
carrying the observation, and a \texttt{predict} function returning a flat
$n_{\text{steps}} \times d_{\text{act}}$ action chunk.

\subsection{Hand-written GGUF reader}
\label{sec:method-reader}

pi0.5 checkpoints are distributed as HuggingFace safetensors. RoboRT
converts them offline to GGUF with two Python scripts
(\texttt{convert\_pi05\_to\_gguf.py} and
\texttt{convert\_pi05\_mmproj\_to\_gguf.py}), then loads them at startup with
a hand-written reader over the ggml \texttt{gguf\_*} API. The reader
supports F32, BF16, and quantized dtypes and maps tensors directly onto the
compute graph without a model-level abstraction layer. The GGUF metadata is
the single source of configuration: the model dimensions, the action
normalization mode (quantile or mean/std), and the denoising step count are
read from the weight file rather than hard-coded into the runtime, so that
retraining or re-quantization never requires recompilation.

\subsection{Cached compute graph}
\label{sec:method-graph}

The inference graph has three shape-dependent inputs: the number of visual
tokens, the number of language tokens, and the suffix length. RoboRT caches
the graph keyed on this triple: the graph is built once per unique shape and
rebuilt only when one of the three changes (e.g., a different language
instruction). All $n_{\text{steps}}$ iterations of one \texttt{predict}
share a single graph, and later \texttt{predict} calls with unchanged shapes
reuse it directly, so the steady-state inference path incurs zero
graph-construction overhead.

\subsection{Flow-matching denoising loop}
\label{sec:method-denoise}

Given a noise tensor $\mathbf{x}_T \sim \mathcal{N}(0, I)$, the action expert
is iterated as
\begin{equation}
  \mathbf{x}_{t-\Delta t} = \mathbf{x}_t - \Delta t \cdot \epsilon_\theta(\mathbf{x}_t, c),
  \label{eq:denoise}
\end{equation}
where $\epsilon_\theta$ is the denoising network conditioned on $c$ (vision,
language, and state) and $\Delta t = 1/n_{\text{steps}}$. The prefix
(vision + language embeddings) is computed once, and its key/value cache is
reused across all denoising iterations. After the loop, the action chunk is
denormalized within the engine using the training statistics (quantile or
mean/std, per the GGUF metadata) and truncated to the real action dimension
$d_{\text{act}}$.

\section{Reproducibility}
\label{sec:repro}

Stochasticity in flow-matching inference originates from the noise sample;
hence, reproducibility reduces to controlling the noise RNG. In RoboRT, the
RNG is maintained as engine state: it advances once per \texttt{predict},
and the environment variable \texttt{VLA\_PI05\_SEED} sets only the starting
point of the sequence. The engine provides \texttt{vla-pi05-selfcheck}, a
deterministic tool that runs with a fixed seed and emits the final
prediction's 1,600 floats. This enables bit-exact comparison across
implementations: any regression in graph construction, operator selection,
or normalization manifests as a numerical discrepancy in the emitted values.
For a faithful comparison, the server must be queried from a clean process;
any additional request advances the noise sequence and alters the marginal
task outcomes.

The serving layer additionally enforces the input contract strictly: the
state length must equal the configured dimension, language token counts must
lie in $[1, n_{\text{lang}}]$, the noise shape must match
$\{0, n_{\text{suffix}} \times d_{\text{act}}\}$, and any float input
containing NaN or Inf is rejected.

\section{Experiments}
\label{sec:experiments}

\subsection{Setup}
We evaluate RoboRT on an A10 GPU (24\,GB) with fp16 weights. The model is
the pi0.5 checkpoint fine-tuned on LIBERO-90. Latency is measured as the
end-to-end \texttt{predict} time with a fixed seed
(\texttt{VLA\_PI05\_SEED=42}). The LeRobot reference processes the same
observation under the same seed in PyTorch.

\subsection{Inference latency}
\label{sec:exp-latency}

Table~\ref{tab:latency} reports the end-to-end latency and its phase
breakdown. RoboRT achieves 161\,ms/step, a 2.3$\times$ speedup over the
LeRobot reference at 377\,ms/step---equivalently 6.2\,Hz versus 2.7\,Hz,
more than halving the gap to the 10\,Hz control target. Within RoboRT,
denoising dominates, accounting for 126.9\,ms (79\% of the total); vision
encoding accounts for 23.7\,ms (15\%), and the remaining serving overhead
amounts to 1.6\,ms. This distribution indicates that further latency
reductions must come primarily from the denoising phase: fewer
flow-matching steps, operator fusion within the denoising layers, and
CUDA-graph capture of the repeated iterations all represent promising
directions for future work.

\begin{table}[t]
  \centering
  \caption{End-to-end latency per inference step on an A10 (fp16).}
  \label{tab:latency}
  \begin{tabular}{lccccc}
    \toprule
    Engine & Vision & Denoise & Other & Total & Speedup \\
    \midrule
    LeRobot (PyTorch) & --- & --- & --- & 377\,ms & 1.0$\times$ \\
    RoboRT (C++)       & 23.7\,ms & 126.9\,ms & 1.6\,ms & \textbf{161\,ms} & 2.3$\times$ \\
    \bottomrule
  \end{tabular}
\end{table}

\subsection{LIBERO evaluation}
\label{sec:exp-libero}

We evaluate the policy on the LIBERO \cite{libero} suite with one episode
per task, a seed offset of 42, and no video recording. RoboRT achieves
\textbf{38/40} task success, and the two failures (tasks \{24, 29\})
coincide exactly with those of the LeRobot reference, confirming that the
engine reproduces the behavior of the reference policy end to end.

\begin{table}[t]
  \centering
  \caption{LIBERO task success (seed offset 42, one episode per task).}
  \label{tab:libero}
  \begin{tabular}{lcc}
    \toprule
    Engine & Success & Failed tasks \\
    \midrule
    LeRobot (reference) & 38/40 & \{24, 29\} \\
    RoboRT (C++)         & 38/40 & \{24, 29\} \\
    \bottomrule
  \end{tabular}
\end{table}

\section{Conclusion}
\label{sec:conclusion}

We have presented RoboRT, a native C++ inference engine for the pi0.5
vision-language-action policy. By matching the runtime to the flow-matching
workload---hand-written GGUF loading, cached compute graphs, and
deterministic noise---RoboRT delivers a 2.3$\times$ latency reduction over
the reference implementation while remaining bit-exact reproducible and
behaviorally identical on LIBERO. The engine is released under the
Apache-2.0 license with conversion scripts, a ZeroMQ/Protobuf serving layer,
and a deterministic self-check tool.

\begin{thebibliography}{9}

\bibitem{pi0}
M.~Black et al.
\newblock $\pi_0$: A vision-language-action flow model for general robot control.
\newblock \emph{arXiv:2410.24164}, 2024.

\bibitem{pi05}
S.~Nasiriany et al.
\newblock $\pi_{0.5}$: A vision-language-action model with open-world generalization.
\newblock \emph{arXiv:2504.09754}, 2025.

\bibitem{lerobot}
R.~Cadene et al.
\newblock LeRobot: State-of-the-art machine learning for real-world robotics.
\newblock \emph{https://github.com/huggingface/lerobot}, 2024.

\bibitem{flowmatching}
Y.~Lipman, R.~T.~Q.~Chen, H.~Ben-Hamu, M.~Nickel, and M.~Le.
\newblock Flow matching for generative modeling.
\newblock \emph{ICLR}, 2023.

\bibitem{ddpm}
J.~Ho, A.~Jain, and P.~Abbeel.
\newblock Denoising diffusion probabilistic models.
\newblock \emph{NeurIPS}, 2020.

\bibitem{sglang}
L.~Zheng et al.
\newblock SGLang: Efficient execution of structured language model programs.
\newblock \emph{arXiv:2312.07104}, 2023.

\bibitem{ggml}
G.~Gerganov et al.
\newblock ggml: Tensor library for machine learning.
\newblock \emph{https://github.com/ggml-org/ggml}, 2023.

\bibitem{siglip}
X.~Zhai, B.~Mustafa, A.~Kolesnikov, and L.~Beyer.
\newblock Sigmoid loss for language image pre-training.
\newblock \emph{ICCV}, 2023.

\bibitem{libero}
B.~Liu, Y.~Zhu, C.~Gao, Y.~Feng, Q.~Liu, Y.~Zhu, and P.~Stone.
\newblock LIBERO: Benchmarking knowledge transfer for lifelong robot learning.
\newblock \emph{NeurIPS}, 2023.

\end{thebibliography}

\end{document}
